\documentclass[11pt]{article}

\usepackage[margin=0.82in]{geometry}
\usepackage[T1]{fontenc}
\usepackage[utf8]{inputenc}
\usepackage{lmodern}
\usepackage{microtype}
\usepackage{xcolor}
\usepackage{hyperref}
\usepackage{booktabs}
\usepackage{longtable}
\usepackage{tabularx}
\usepackage{array}
\usepackage{enumitem}
\usepackage{listings}
\usepackage{amsmath,amssymb}
\usepackage{fancyhdr}
\usepackage{titlesec}
\usepackage{lastpage}

\definecolor{navy}{HTML}{16324F}
\definecolor{teal}{HTML}{167D7F}
\definecolor{orange}{HTML}{D97706}
\definecolor{lightblue}{HTML}{EAF3F8}
\definecolor{lightgray}{HTML}{F4F5F7}
\definecolor{codegreen}{HTML}{276749}
\definecolor{codered}{HTML}{9B2C2C}

\hypersetup{
  colorlinks=true,
  linkcolor=navy,
  urlcolor=teal,
  citecolor=orange,
  pdftitle={Wildfire Modeling and Machine Learning Research Roadmap},
  pdfauthor={Research onboarding guide}
}

\titleformat{\section}{\Large\bfseries\color{navy}}{\thesection}{0.7em}{}
\titleformat{\subsection}{\large\bfseries\color{teal}}{\thesubsection}{0.7em}{}
\setlist{nosep,leftmargin=1.5em}
\setlength{\parindent}{0pt}
\setlength{\parskip}{0.55em}
\renewcommand{\arraystretch}{1.18}

\pagestyle{fancy}
\fancyhf{}
\lhead{Wildfire ML Research Roadmap}
\rhead{\thepage\ of \pageref{LastPage}}
\renewcommand{\headrulewidth}{0.3pt}

\lstdefinestyle{terminal}{
  backgroundcolor=\color{lightgray},
  basicstyle=\ttfamily\small,
  breaklines=true,
  frame=single,
  rulecolor=\color{navy!35},
  keywordstyle=\color{navy}\bfseries,
  commentstyle=\color{codegreen},
  stringstyle=\color{codered},
  showstringspaces=false,
  columns=fullflexible,
  keepspaces=true
}
\lstset{style=terminal}

\newcommand{\important}[1]{%
  \begin{center}
  \fcolorbox{orange}{orange!8}{\parbox{0.93\linewidth}{\textbf{Important:} #1}}
  \end{center}}

\newcommand{\deliverable}[1]{\textcolor{teal}{\textbf{Deliverable:}} #1}

\begin{document}

\begin{titlepage}
\centering
\vspace*{1.2cm}
{\Huge\bfseries\color{navy} Wildfire Modeling and Machine Learning\par}
\vspace{0.25cm}
{\LARGE\color{teal} Research Onboarding Roadmap\par}
\vspace{0.8cm}
\rule{0.85\textwidth}{1pt}
\vspace{0.8cm}

{\large A practical guide to WRF-SFIRE, WRF-Chem, FARSITE/FlamMap,\par
research objectives, experiments, evaluation, and career preparation\par}

\vfill
\begin{tabular}{rl}
Student: & \textit{Your name} \\
Advisor: & \textit{Professor's name} \\
Laboratory: & \textit{Research group} \\
Prepared: & October 3, 2026 \\
Status: & Living research plan; confirm scope with advisor
\end{tabular}
\vfill

\important{The project description below is a reasoned interpretation of the first meeting, not a substitute for the advisor's exact research question. The first research task is to convert the professor's idea into a one-page, measurable problem statement.}
\end{titlepage}

\tableofcontents
\newpage

\section{Executive answer: what this project probably is}

The most defensible interpretation is:

\begin{quote}
Build and rigorously evaluate a computationally cheaper machine-learning surrogate (or hybrid physics--ML system) for a carefully selected part of WRF-SFIRE, so that it predicts fire evolution with accuracy close to the physics baseline while using much less wall time, memory, and energy.
\end{quote}

WRF-SFIRE is \emph{not} itself an AI model. It is a coupled numerical weather--wildfire model. WRF computes the evolving atmosphere; SFIRE advances a fire front, commonly with a level-set method and semi-empirical spread-rate formulas. Wind drives the fire, and fire heat and moisture fluxes feed back into the atmosphere \cite{mandel2011}. This two-way coupling is scientifically valuable and computationally expensive.

The phrase ``predict fire'' is ambiguous. It could mean:
\begin{enumerate}
  \item \textbf{ignition occurrence}: will a fire start in a place and time?
  \item \textbf{fire spread}: where will the perimeter be over the next hours?
  \item \textbf{fire behavior}: rate of spread, flame length, intensity, or heat flux;
  \item \textbf{smoke and air quality}: where will emissions and particulate matter travel?
  \item \textbf{risk}: burn probability or expected impact under many scenarios.
\end{enumerate}

The repositories and software mentioned in the meeting point most strongly to items 2--4. Do not train a model until the advisor selects the target.

\important{``Same or better accuracy with less GPU'' needs correction. Traditional WRF-SFIRE is primarily a CPU-parallel MPI/OpenMP scientific code; a GPU is not normally required for the basic build. The proposed ML system may use a GPU for training and inference. The comparison must therefore report wall-clock time, CPU/GPU hours, peak RAM/VRAM, energy if possible, and predictive error.}

\section{What each repository or program does}

\begin{longtable}{p{0.19\textwidth}p{0.29\textwidth}p{0.23\textwidth}p{0.21\textwidth}}
\toprule
\textbf{Tool} & \textbf{Purpose} & \textbf{Inputs/outputs} & \textbf{Likely role in your project} \\
\midrule
\endfirsthead
\toprule
\textbf{Tool} & \textbf{Purpose} & \textbf{Inputs/outputs} & \textbf{Likely role in your project} \\
\midrule
\endhead
WRF-SFIRE & Two-way coupled atmosphere and wildland-fire simulation built on WRF. Current repository documentation labels the release W4.4-S0.1 and selects SFIRE with \texttt{ifire=1} \cite{wrfsfireRepo}. & Terrain, fuels, moisture, atmospheric initial/boundary conditions, ignition; outputs gridded weather, fire arrival/spread, heat flux, and related fields. & High-fidelity teacher/baseline and synthetic-data generator. It may also be retained as a physics component in a hybrid system. \\
\addlinespace
WRF-Chem & WRF coupled with atmospheric chemistry. It simulates gases, aerosols, emissions, transport, chemical reactions, and deposition. NCAR states that WRF-Chem is no longer under active development, although selected schemes and preprocessing tools remain supported \cite{wrfchem}. & Meteorology, emissions, chemistry initial/boundary conditions; outputs 3-D chemical and aerosol fields. & Relevant if the project includes smoke, PM$_{2.5}$, or air quality. It is not necessary for a first fire-perimeter surrogate. \\
\addlinespace
WRF-Chem Preprocessing Tools & NCAR utilities that prepare chemical boundary/initial conditions and anthropogenic, biogenic, aircraft, and fire-emission inputs. The \texttt{fire\_emiss} tool processes FINN fire-emission inventory data \cite{chemtools}. & External inventories/global models to WRF-Chem-ready NetCDF files. & Data preparation for a smoke/chemistry branch, not a substitute for WRF-SFIRE and not an ML repository. \\
\addlinespace
FlamMap 6 & Windows desktop fire-analysis system. It can compute potential fire behavior, spread, conditional burn probabilities, and longer fire growth under specified fuels, terrain, moisture, wind, and weather \cite{flammap}. & An eight-layer landscape file (elevation, slope, aspect, fuel model, canopy cover, height, base height, and bulk density), plus weather/wind/moisture; outputs spread rate, flame length, intensity, crown-fire activity, growth, burn probability, emissions, and consumption. & Independent baseline, domain-learning tool, or source of scenario data. Easier first exposure than a real-data WRF-SFIRE workflow, but it is a different model. \\
\addlinespace
FARSITE & Fire Area Simulator for time-varying fire growth. The Forest Service states that FARSITE 4 is no longer supported or downloadable and is now included in FlamMap 6 \cite{farsite}. & Landscape, fuels, weather, wind, moisture, ignition; outputs time-evolving fire growth and behavior. & A comparison baseline through FlamMap 6. Do not search for a separate modern FARSITE 4 installer. \\
\addlinespace
\texttt{far.site} & Commercial Federal Acquisition Regulation compliance software, unrelated to wildfire modeling \cite{farsiteWrong}. & Government-contract compliance data. & None. This link was probably a transcription or naming mix-up; ask whether the professor meant FARSITE. \\
\bottomrule
\end{longtable}

\subsection{What the NASA connection does and does not mean}

NASA Applied Sciences has described WRF-SFIRE as an open-source system for forecasts of fuel moisture, fire spread, smoke, and weather, and NASA-supported work has used it in fire applications \cite{nasaFire}. This supports the claim that the system is credible and relevant to NASA-funded or NASA-used research. It does \emph{not}, by itself, prove a formal ``NASA certification'' or that the GitHub repository passed a universal NASA acceptance test. In a paper or presentation, use the precise wording ``used/described in NASA Applied Sciences work'' unless the professor supplies a specific program, award, benchmark, or acceptance document.

\section{A strong project objective}

\subsection{Recommended first objective}

Start narrower than a complete replacement of WRF-SFIRE:

\begin{quote}
Given terrain, fuel category, fuel moisture, ignition geometry, and near-surface wind fields, predict the fire-arrival-time field (or a sequence of burned-area masks) over a fixed forecast horizon. Compare against held-out WRF-SFIRE simulations and simple baselines. Demonstrate a pre-agreed error tolerance and a pre-agreed speed/memory improvement on the same hardware.
\end{quote}

This objective is achievable because the output is spatial, physically interpretable, and convertible into perimeters at any forecast time. Fire arrival time $T(x,y)$ yields a burned mask at time $t$ by
\begin{equation}
B_t(x,y)=\mathbb{1}\{T(x,y)\le t\}.
\end{equation}

After this succeeds, extend to uncertainty, real terrain, changing weather, coupled atmospheric feedback, or smoke.

\subsection{Three possible scopes to confirm with the advisor}

\begin{enumerate}
  \item \textbf{Level A: fire-spread surrogate.} Prescribed wind/terrain/fuels $\rightarrow$ arrival time or perimeter. Lowest complexity and best starting point.
  \item \textbf{Level B: coupled fire--weather emulator.} Initial atmospheric and land state $\rightarrow$ both fire and selected near-fire weather variables. More faithful, much harder, and autoregressive stability matters.
  \item \textbf{Level C: fire--smoke/chemistry system.} Fire and weather $\rightarrow$ emission injection and downwind smoke/PM fields. Requires WRF-Chem knowledge and preprocessing; it should be a later milestone unless this is explicitly the lab's central goal.
\end{enumerate}

\subsection{Proposed success criteria}

Replace the bracketed values after discussion:
\begin{itemize}
  \item median fire-arrival-time MAE $\le$ [10--20] minutes over burned cells;
  \item perimeter/burned-area IoU $\ge$ [0.80] at [1, 2, 4, 6] hours;
  \item 95th-percentile symmetric boundary distance $\le$ [two fire-grid cells];
  \item final burned-area relative error $\le$ [10\%];
  \item inference speedup $\ge$ [50--100$\times$] versus the defined WRF-SFIRE workflow;
  \item peak memory reduction $\ge$ [value] or inference within [hardware limit];
  \item no physically impossible unburning, no spread through non-burnable cells, and stable rollouts;
  \item acceptable performance on geographically and meteorologically out-of-distribution tests.
\end{itemize}

Accuracy and efficiency are a Pareto tradeoff. Report the frontier rather than claiming one model is universally best.

\section{Questions to take to the next advisor meeting}

Ask these before spending weeks installing software or training models:

\begin{enumerate}
  \item What exact target is intended: ignition, perimeter, arrival time, intensity, smoke, or air quality?
  \item Does ``same accuracy'' mean agreement with WRF-SFIRE, agreement with observed fires, or both?
  \item Which geographic region, spatial resolution, forecast horizon, and update frequency matter?
  \item Is the desired model a full surrogate, a correction to a physics model, a downscaler, or a data-assimilation component?
  \item Is there an existing simulation/observation dataset, configuration, benchmark case, or prior student's code?
  \item Which WRF-SFIRE commit/tag, namelists, WPS data, and compiler environment are the lab standard?
  \item What did ``accepted by NASA'' refer to: funding, deployment, validation, publication, or a specific benchmark?
  \item Did the professor mean FARSITE rather than \url{https://far.site/}?
  \item What compute is available: workstation specifications, university HPC allocation, cloud credits, storage, and job limits?
  \item What is the expected first deliverable and date: literature review, reproducible run, dataset, baseline model, paper, or demo?
  \item What restrictions apply to operational or observational data, and who owns resulting code and data?
  \item Which publication venue or end-user decision is the project designed to support?
\end{enumerate}

\section{How to build and run WRF-SFIRE}

\subsection{Expectations before installation}

The repository is a large Fortran/C scientific application, not a Python package. Compilation is normal. The idealized \texttt{em\_fire} case avoids WPS and real weather data, so it is the correct first milestone. A real-fire case adds WPS, static geographic data, meteorological GRIB files, fuel/terrain preparation, nesting, boundary conditions, and much more storage.

The official WRF guide identifies NetCDF as mandatory. MPI is needed for distributed-memory runs. WPS handling of GRIB2 normally needs zlib, libpng, and JasPer. Compiler/library consistency matters \cite{wrfCompile}. The older OpenWFM compile page gives the project-specific sequence \texttt{configure}, \texttt{compile em\_fire}, and optionally \texttt{compile em\_real} \cite{openwfmCompile}.

\subsection{Recommended environment}

Use Linux. On Windows, use a native Linux workstation, university cluster, or WSL2 with enough RAM and disk; FlamMap itself requires 64-bit Windows \cite{flammap}. Record all versions. For reproducible research, a lab-maintained container or HPC module recipe is better than undocumented manual setup.

Suggested minimum for an idealized learning run: 4--8 modern CPU cores, 16 GB RAM, and 30--60 GB free disk. A real high-resolution coupled case may need far more cores, RAM, and storage. These are planning estimates, not repository guarantees.

\subsection{Ubuntu/Debian quick-start for the idealized fire case}

Commands below are a practical GNU/OpenMPI route. Package names vary by operating system. Do not blindly run them on a managed cluster; use its modules and ask the administrator.

\begin{lstlisting}[language=bash,caption={Install a consistent GNU toolchain and common dependencies}]
sudo apt update
sudo apt install -y build-essential gfortran m4 csh perl git \
  libnetcdf-dev libnetcdff-dev netcdf-bin \
  openmpi-bin libopenmpi-dev

gcc --version
gfortran --version
nc-config --all
nf-config --all
mpif90 --version
\end{lstlisting}

Clone a fixed revision. For initial exploration, \texttt{master} is convenient; for research, pin the exact commit and never let a moving branch silently change the dataset.

\begin{lstlisting}[language=bash,caption={Clone and record provenance}]
git clone https://github.com/openwfm/WRF-SFIRE.git
cd WRF-SFIRE
git checkout master
git rev-parse HEAD | tee SOURCE_COMMIT.txt
git submodule update --init --recursive
\end{lstlisting}

On typical Ubuntu packages, NetCDF is under \texttt{/usr}. Confirm with \texttt{nc-config} and \texttt{nf-config}. Then configure:

\begin{lstlisting}[language=bash,caption={Configure WRF-SFIRE}]
export NETCDF=/usr
export CC=gcc
export CXX=g++
export FC=gfortran
export F77=gfortran
export F90=gfortran
export OMP_NUM_THREADS=1

./clean -a
./configure
\end{lstlisting}

Choose a GNU option that actually exists in the displayed menu. For a first laptop test, serial is easiest; for useful multicore execution, choose GNU \texttt{dmpar} and basic nesting (option 1). Save the terminal choice in a lab notebook.

Compile the ideal case. If parallel compilation hides the first error, retry with \texttt{-j 1}.

\begin{lstlisting}[language=bash,caption={Compile and verify}]
./compile em_fire -j 4 >& compile_em_fire.log
tail -n 40 compile_em_fire.log
grep -n -i "error" compile_em_fire.log | head
ls -lh main/ideal.exe main/wrf.exe
\end{lstlisting}

The repository states that compilation creates links and test subdirectories under \texttt{test/em\_fire}; the default linked case is \texttt{hill\_simple} \cite{emFireReadme}. Initialize and run it:

\begin{lstlisting}[language=bash,caption={Run the idealized hill case}]
cd test/em_fire
grep -n "ifire" namelist.input
./ideal.exe >& ideal.log

# Serial build:
./wrf.exe >& wrf.log

# For a dmpar build, use this instead of the serial command:
# mpirun -np 4 ./wrf.exe >& wrf.log

tail -n 30 wrf.log
grep -n "SUCCESS COMPLETE WRF" rsl.error.0000 wrf.log 2>/dev/null
ls -lh wrfinput_d01 wrfout_d01_*
\end{lstlisting}

Check that \texttt{ifire=1} selects WRF-SFIRE. Do not assume a run is correct merely because an output file exists: inspect the log, timestamps, fields, and spatial evolution. Use Python packages such as \texttt{xarray}, \texttt{netCDF4}, \texttt{wrf-python}, \texttt{matplotlib}, and \texttt{cartopy} for analysis, preferably in a separate environment.

\subsection{Troubleshooting order}

\begin{enumerate}
  \item Read the \emph{first} compiler error, not the last cascading error.
  \item Verify \texttt{gcc}, \texttt{gfortran}, NetCDF-C, NetCDF-Fortran, and MPI use compatible toolchains.
  \item Confirm \texttt{nc-config --prefix}, \texttt{nf-config --prefix}, headers, and libraries agree.
  \item Run \texttt{./clean -a}, reconfigure, and compile with one job.
  \item Check disk space, memory, stack limits, and permissions.
  \item On MPI failure, first prove a serial build/run, then isolate the MPI configuration.
  \item Preserve \texttt{configure.wrf}, compile logs, namelists, commit SHA, environment/module list, and commands with every experiment.
\end{enumerate}

\subsection{Moving from idealized to real-data simulations}

A real case generally requires:
\begin{enumerate}
  \item compile \texttt{em\_real};
  \item compile a compatible WPS version against the WRF-SFIRE build;
  \item download WPS static geography data and meteorological forcing (for example, an appropriate analysis/forecast product);
  \item configure domains and nests in \texttt{namelist.wps};
  \item run \texttt{geogrid.exe}, link GRIB files, run \texttt{ungrib.exe}, and run \texttt{metgrid.exe};
  \item prepare fire-grid terrain, fuels, moisture, and ignition consistently with the SFIRE namelist;
  \item run \texttt{real.exe}, validate boundary/initial fields, then run \texttt{wrf.exe};
  \item compare against observations only after checking coordinate systems, timestamps, units, and spin-up.
\end{enumerate}

This is where the professor's known-good case is extremely valuable. Reproducing a lab case is safer than inventing one from scratch.

\section{A disciplined ML research workflow}

\subsection{Phase 0: define the claim before the model}

Write one page containing:
\begin{itemize}
  \item input fields and their units/resolution;
  \item target fields and forecast horizons;
  \item operational constraints and target hardware;
  \item in-distribution and out-of-distribution test definitions;
  \item accuracy, speed, memory, and uncertainty metrics;
  \item baselines and ablations;
  \item what would falsify the hypothesis.
\end{itemize}

\deliverable{Advisor-approved problem statement and evaluation protocol.}

\subsection{Phase 1: reproduce the physics baseline}

Run one idealized case unchanged. Then vary only one factor at a time (wind speed/direction, fuel, moisture, slope, or ignition). Confirm that outputs respond physically. Profile runtime, CPU utilization, memory, output size, and failure modes.

\deliverable{A script/notebook producing plots, metrics, logs, and a reproducibility manifest from a clean run.}

\subsection{Phase 2: design a dataset, not just a pile of runs}

Construct a parameter table before simulation. Use space-filling or stratified sampling across plausible ranges, and include edge cases. Store:
\begin{itemize}
  \item static inputs: elevation, slope, aspect, fuel type/properties, canopy fields;
  \item dynamic inputs: wind components, temperature, humidity, fuel moisture, pressure, and possibly boundary forcing;
  \item ignition: location, time, geometry, and uncertainty;
  \item targets: fire arrival time, burned mask sequence, spread rate, heat flux, and selected atmospheric fields;
  \item provenance: source commit, configuration, seed, compiler, machine, start/end time, checksums, and run status.
\end{itemize}

Use chunked formats such as Zarr or NetCDF/HDF5 with explicit metadata. Keep raw outputs immutable; create versioned processed datasets. Build automatic quality-control tests for NaNs, units, monotonic burn progression, non-burnable masks, time ordering, and coordinate alignment.

Avoid leakage. Randomly splitting neighboring frames from the same simulation makes test results unrealistically easy. Split by complete simulation, event, geography, and preferably weather regime. Reserve a truly untouched final test set.

\deliverable{Versioned dataset card, generation manifest, QC report, and train/validation/test split file.}

\subsection{Phase 3: establish baselines}

Before a sophisticated network, evaluate:
\begin{enumerate}
  \item persistence/no-spread and simple radial/constant-rate predictors;
  \item a reduced fire-spread calculation if available;
  \item a small convolutional encoder--decoder or U-Net;
  \item only then, operator-learning or sequence architectures such as FNO, U-Net rollouts, ConvLSTM, or a transformer if justified.
\end{enumerate}

The winning method is the smallest one that meets the agreed accuracy and deployment constraints. Parameter count alone is not compute cost; report measured latency and memory.

\subsection{Phase 4: add physics and constraints}

Useful constraints may include:
\begin{itemize}
  \item burned area cannot become unburned;
  \item no spread through non-burnable fuels unless spotting is explicitly modeled;
  \item nonnegative arrival-time increments and heat release;
  \item rotation/reflection handling only when terrain, wind, and coordinates are transformed consistently;
  \item losses emphasizing the moving boundary rather than overwhelming background pixels;
  \item conservation or budget checks for any predicted fluxes;
  \item uncertainty estimates or ensembles because a single deterministic perimeter can be dangerously overconfident.
\end{itemize}

A hybrid model may be scientifically stronger than a black-box replacement: keep a known spread law or coarse solver and learn a correction, closure, downscaling, or unresolved coupling term.

\subsection{Phase 5: evaluate like an environmental model}

Use several metric families:

\begin{longtable}{p{0.22\textwidth}p{0.69\textwidth}}
\toprule
\textbf{Question} & \textbf{Metrics/checks} \\
\midrule
Arrival timing & MAE/RMSE/bias of $T(x,y)$ over valid/burned regions; stratify by horizon and rate-of-spread regime. \\
Spatial overlap & IoU/Jaccard, Dice/F1, precision, recall, false alarm ratio, probability of detection. \\
Boundary location & Hausdorff distance (prefer robust percentile), average symmetric surface distance, perimeter displacement. \\
Size/shape & Absolute and relative burned-area error, perimeter length error, spread-direction error, connected components. \\
Probabilistic skill & Reliability diagram, Brier score, CRPS, interval coverage and sharpness. \\
Physical behavior & Monotonicity, barrier violations, unrealistic islands, stability over long rollouts, sensitivity to wind/slope/moisture. \\
Efficiency & End-to-end latency including preprocessing and I/O, throughput, CPU/GPU hours, peak RAM/VRAM, energy, model/dataset storage. \\
Robustness & New ignitions, winds, fuels, topography, grid spacing, regions, seasons, and perturbations; report failure cases. \\
\bottomrule
\end{longtable}

Compare on the same machine or normalize carefully. Include WRF-SFIRE initialization and required preprocessing if the claimed use case includes them; do not compare only a neural forward pass with an entire physics workflow without explaining the boundary.

\subsection{Phase 6: validate against observations}

Agreement with WRF-SFIRE measures imitation, not truth. A publishable system should eventually compare with independent observations such as quality-controlled fire perimeters/arrival estimates, weather stations, airborne data, or satellite products. Satellite active-fire detections have missing observations, clouds, resolution limits, and timing uncertainty; treat them statistically, not as perfect labels \cite{mandelAssimilation}.

Separate error sources: atmospheric forcing, fuels, moisture, ignition, physics, numerical resolution, observation error, and surrogate error. Otherwise ``ML error'' is not interpretable.

\section{Suggested 12-week onboarding plan}

\begin{longtable}{p{0.10\textwidth}p{0.54\textwidth}p{0.28\textwidth}}
\toprule
\textbf{Week} & \textbf{Work} & \textbf{Evidence of completion} \\
\midrule
1 & Confirm scope using the advisor questions; read the WRF-SFIRE model paper and repository README. & One-page objective and glossary. \\
2 & Prepare Linux/toolchain; compile \texttt{em\_fire}; freeze commit and environment. & Executables, logs, manifest. \\
3 & Run and visualize the unchanged idealized case; identify variables, dimensions, units. & Reproducible notebook and figures. \\
4 & Sensitivity experiments varying one input at a time; profile runtime and memory. & Physics sanity-check report. \\
5 & Finalize parameter ranges, storage schema, QC rules, and split policy. & Dataset design review. \\
6 & Generate a small pilot ensemble; automatically detect failed or corrupted runs. & Pilot dataset v0.1 and QC dashboard. \\
7 & Implement persistence/simple spread and small CNN/U-Net baselines. & Baseline table with fixed test cases. \\
8 & Implement the main candidate model and training pipeline with seeds/checkpoints. & Reproducible training run. \\
9 & Add boundary-aware/physics constraints and conduct ablations. & Ablation table. \\
10 & Test new winds, fuels, slopes, and ignitions; measure calibration and rollout stability. & Robustness/error atlas. \\
11 & Benchmark end-to-end CPU/GPU time, memory, energy if available, and accuracy frontier. & Pareto plot and profiling report. \\
12 & Present results, limitations, failure cases, and next-stage real-data plan. & Advisor presentation and research memo. \\
\bottomrule
\end{longtable}

If compilation or data access consumes more than a week, ask for the lab's known-good environment rather than silently losing a month.

\section{Repository and experiment structure}

Do not mix generated terabytes, upstream WRF source, and ML code in one unstructured directory. A reasonable layout is:

\begin{lstlisting}[caption={Suggested project organization}]
wildfire-surrogate/
  README.md
  CITATION.cff
  environment/          # lock files, container recipe, module list
  configs/              # versioned YAML/TOML experiment configs
  scripts/              # dataset generation, launch, evaluation
  src/                  # reusable Python package
  tests/                # unit and physics-invariant tests
  notebooks/            # exploration only, not the sole pipeline
  metadata/             # manifests, split files, checksums
  reports/              # figures, tables, manuscript sources
  data/README.md         # locations and policy, not giant files in Git
  wrf-sfire/             # pinned upstream clone/submodule if approved
\end{lstlisting}

Use Git for code/configuration, not giant NetCDF outputs. Use institutional storage, object storage, or a data registry for generated data. Every result should be reconstructable from a commit, configuration, dataset version, seed, and environment.

\section{Common failure modes and how to avoid them}

\begin{itemize}
  \item \textbf{Undefined target:} ``predict fire'' produces the wrong dataset. Fix the scientific question first.
  \item \textbf{Teacher equals truth:} matching WRF-SFIRE only proves emulation. Validate observations separately.
  \item \textbf{Frame leakage:} nearby times from one run appear in train and test. Split at event/simulation level.
  \item \textbf{Background-dominated loss:} predicting ``unburned everywhere'' looks accurate. Use boundary and class-aware metrics.
  \item \textbf{No easy baseline:} a complex network beats nothing. Include persistence and simple physics/statistical baselines.
  \item \textbf{Unfair speed claim:} CPU physics model versus GPU neural kernel with different I/O boundaries. Define end-to-end timing and hardware.
  \item \textbf{Overfitting idealized cases:} strong test scores collapse on real terrain/weather. Design distribution shifts from day one.
  \item \textbf{Ignoring uncertainty:} a crisp perimeter implies false confidence. Calibrate probabilities or ensembles.
  \item \textbf{Changing upstream version:} moving \texttt{master} changes labels. Pin SHA/tag and archive configs.
  \item \textbf{Running everything on a laptop:} start with idealized smoke tests locally; move ensemble generation to HPC.
  \item \textbf{Starting with WRF-Chem:} chemistry greatly expands state and cost. Add it only if smoke/air quality is the defined target.
\end{itemize}

\section{Skills and prerequisites to study}

\subsection{Scientific foundation}

Learn enough to detect nonsense:
\begin{itemize}
  \item wildfire behavior: fuels, fuel moisture, Rothermel-type spread, slope/wind effects, crown fire, spotting limitations;
  \item atmospheric science: wind fields, boundary layer, stability, humidity, terrain flows, nesting, initial/boundary conditions;
  \item numerical modeling: grids, CFL stability, discretization, level-set methods, parameterization, convergence;
  \item geospatial science: projections, raster alignment, resampling, GeoTIFF/NetCDF, GIS, LANDFIRE-style landscape layers;
  \item uncertainty and validation: observational error, calibration, sensitivity, ensembles, data assimilation.
\end{itemize}

\subsection{Engineering and ML foundation}

Prioritize Linux/bash, Git, Python, NumPy, xarray, NetCDF/Zarr, matplotlib/cartopy, PyTorch or JAX, configuration management, testing, profiling, Slurm, MPI concepts, and basic Fortran reading. Learn Docker/Apptainer after the first manual build so the environment becomes reproducible.

For ML, master convolutional segmentation, spatiotemporal forecasting, class imbalance, autoregressive error, neural operators, uncertainty calibration, and distribution shift. The architecture name is less important than data design and evaluation.

\section{Career value}

Yes, this can help significantly with a job if the outcome is more than ``I trained a neural network.'' The combination is uncommon and valuable:

\begin{itemize}
  \item scientific ML and physics-informed modeling;
  \item high-performance computing, MPI workflows, and profiling;
  \item geospatial and Earth-observation data engineering;
  \item environmental risk, weather, wildfire, smoke, and climate applications;
  \item reproducible research and uncertainty-aware evaluation.
\end{itemize}

Relevant roles include scientific ML engineer, research software engineer, geospatial ML engineer, environmental data scientist, weather/climate modeler, remote-sensing scientist, HPC engineer, and applied scientist in government laboratories, universities, weather/climate companies, insurance/catastrophe modeling, utilities, emergency management, Earth-observation, and sustainability teams.

The strongest portfolio artifacts would be:
\begin{enumerate}
  \item a reproducible WRF-SFIRE baseline with documented environment;
  \item a public or shareable dataset card and generation/QC pipeline;
  \item an honest benchmark including accuracy, uncertainty, latency, memory, and failure cases;
  \item a clean tested repository with a small demo;
  \item a poster, workshop paper, preprint, or peer-reviewed publication;
  \item evidence of collaboration with domain scientists and clear communication of limitations.
\end{enumerate}

On a resume, quantify the contribution: for example, ``Built a reproducible surrogate for fire-arrival-time prediction on held-out simulations; achieved X IoU and Y-minute arrival-time MAE with Z-fold end-to-end speedup and Q\% lower peak memory.'' Never fill those letters with training-set results.

\section{The first five actions to take now}

\begin{enumerate}
  \item Send the advisor a concise summary: target, proposed first objective, uncertainties, and the 12 questions in Section 5.
  \item Obtain the lab's exact WRF-SFIRE version, known-good test/real case, and computing environment.
  \item Compile and run only \texttt{em\_fire}; capture the commit, logs, namelist, outputs, runtime, and memory.
  \item Read and plot the fire-arrival/perimeter-related variables; perform one wind and one slope sensitivity test.
  \item Present a one-page dataset/evaluation design before launching a large simulation ensemble.
\end{enumerate}

\section{A concise meeting-ready project pitch}

\begin{quote}
I propose to begin with a reproducible idealized WRF-SFIRE benchmark and learn a surrogate for fire arrival time from terrain, fuels, moisture, ignition, and wind. I will split by complete scenarios, evaluate arrival-time and perimeter errors plus physical constraints, and measure end-to-end speed, RAM/VRAM, and energy on fixed hardware. I will first prove generalization to unseen simulated conditions, then validate selected cases against observations. After that, we can decide whether to extend to two-way atmospheric feedback or WRF-Chem smoke prediction.
\end{quote}

\section{Source notes and links}

The primary sources below should be rechecked when beginning the work because repositories, documentation, and downloads change. Access dates for web pages: October 3, 2026.

\begin{thebibliography}{99}

\bibitem{mandel2011}
J. Mandel, J. D. Beezley, and A. K. Kochanski,
``Coupled atmosphere-wildland fire modeling with WRF 3.3 and SFIRE 2011,''
\emph{Geoscientific Model Development}, 4, 591--610, 2011.
\url{https://doi.org/10.5194/gmd-4-591-2011}

\bibitem{wrfsfireRepo}
OpenWFM, ``WRF-SFIRE source repository and README.''
\url{https://github.com/openwfm/WRF-SFIRE}

\bibitem{emFireReadme}
OpenWFM, ``WRF-SFIRE idealized fire test README.''
\url{https://github.com/openwfm/WRF-SFIRE/blob/master/test/em_fire/README.txt}

\bibitem{openwfmCompile}
OpenWFM Wiki, ``How to compile WRF-SFIRE.''
\url{https://wiki.openwfm.org/wiki/How_to_compile_WRF-Fire}

\bibitem{wrfCompile}
WRF Model, ``WRF Users Guide: Compiling WRF.''
\url{https://github.com/wrf-model/Users_Guide/blob/main/compiling.rst}

\bibitem{wrfchem}
NSF NCAR Atmospheric Chemistry Observations \& Modeling,
``WRF-Chem.''
\url{https://www2.acom.ucar.edu/wrf-chem}

\bibitem{chemtools}
NSF NCAR, ``WRF-Chem Preprocessing Tools.''
\url{https://github.com/NCAR/WRF-Chem-Preprocessing-Tools}

\bibitem{flammap}
USDA Forest Service, Missoula Fire Sciences Laboratory, ``FlamMap.''
\url{https://research.fs.usda.gov/firelab/products/dataandtools/flammap}

\bibitem{farsite}
USDA Forest Service, Missoula Fire Sciences Laboratory, ``FARSITE.''
\url{https://research.fs.usda.gov/firelab/products/dataandtools/farsite}

\bibitem{farsiteWrong}
Control Insights, ``FAR SITE: Federal Acquisition Regulation compliance software.''
\url{https://far.site/}

\bibitem{nasaFire}
NASA Applied Sciences Disasters Program, ``Fires'' (WRF-SFIRE description and applications).
\url{https://disasters.nasa.gov/what-we-do/disasters/fires}

\bibitem{mandelAssimilation}
J. Mandel, A. K. Kochanski, M. Vejmelka, and J. D. Beezley,
``Data assimilation of satellite fire detection in coupled atmosphere-fire simulation by WRF-SFIRE,'' 2014.
\url{https://arxiv.org/abs/1410.6948}

\end{thebibliography}

\appendix
\section{Advisor confirmation form}

\begin{tabularx}{\textwidth}{>{\bfseries}p{0.31\textwidth}X}
Primary prediction target & \hrulefill \\
Forecast horizon/update rate & \hrulefill \\
Region and resolution & \hrulefill \\
Teacher/baseline version & \hrulefill \\
Observational validation data & \hrulefill \\
Available compute/storage & \hrulefill \\
Accuracy threshold & \hrulefill \\
Speed/memory threshold & \hrulefill \\
First deliverable and due date & \hrulefill \\
Expected publication/end user & \hrulefill \\
\end{tabularx}

\section{Reproducibility checklist}

For every reported result, preserve:
\begin{itemize}
  \item code repository URL and exact commit SHA;
  \item upstream WRF-SFIRE commit/tag and local patches;
  \item compiler, MPI, NetCDF, operating system, CPU/GPU, and library versions;
  \item complete namelists/configuration and random seeds;
  \item raw-data origin, license, checksums, processing version, and split IDs;
  \item training logs, checkpoints, optimizer/schedule, stopping rule, and selected epoch;
  \item exact metric code and inclusion/exclusion masks;
  \item timing protocol, warm-up policy, repeats, batch size, and I/O boundary;
  \item failed runs and excluded samples with reasons;
  \item figures/tables generated automatically from archived outputs.
\end{itemize}

\section{Glossary}

\begin{description}[leftmargin=2.8cm,style=nextline]
  \item[WRF] Weather Research and Forecasting model, a numerical atmospheric model.
  \item[SFIRE] Spread FIRE component coupled to WRF.
  \item[WPS] WRF Preprocessing System for domains, geographic data, and meteorological inputs.
  \item[WRF-Chem] WRF coupled with atmospheric chemistry and aerosols.
  \item[FINN] Fire INventory from NCAR, used for fire emissions.
  \item[NetCDF] Self-describing array-data format widely used for atmospheric model input/output.
  \item[MPI] Message Passing Interface for distributed-memory parallel execution.
  \item[Level set] Numerical representation of a moving boundary as a contour of a scalar field.
  \item[Surrogate] A learned approximation intended to reproduce selected input--output behavior of a costly model.
  \item[Emulator] Often used synonymously with surrogate; may additionally model uncertainty.
  \item[IoU] Intersection over union of predicted and reference burned regions.
  \item[OOD] Out of distribution: conditions meaningfully different from training data.
  \item[Data assimilation] Statistical combination of observations with model state/forecast.
\end{description}

\end{document}
